Fragestellungen Ideen:
(1) "Wie viel Privatsphäre würdest du aufgeben"
-> was wäre für mich ok?" - in schöner formuliert :D

Extremwerte: 
- der Zweck heiligt die Mittel (ich würde meine Privatsphäre zum Schutz aufgeben)
- Privatsphäre ist mir heilig (niemand darf meine Daten haben, ich nutze kein Social Media)

(2) Wo zieht man die Grenze?
Welche Mittel sind noch ok und welche nicht mehr?
- unter Verdacht für die meisten ok

(3) false-positives - sind die ok?
-> relativ eindeutige Anwort :(

(x) Alternative Umsetzungen




Sowohl Präsentation/Diskussion als auch die Hausarbeit sollte behandeln, dass das ganze keine schwarz/weiß-Frage ist, sondern es sich um ein "Spektrum" (doofes Wort dafür) handelt. In der Diskussion kommt das durch die Einordnung in den ersten zwei Fragen. In der Hausarbeit kommt das durch die ganze Diskussion und es kann super als Grundlage für's Fazit verwendet werden.

- danke :D


Gliederung:
- Schutz der Kinder: CSA-Verordnung Grundgedanke
- Chatkontrolle: E2E \& Client-Side-Scanning
- Analyse durch KI
-> Analyse unstrukturierter Daten: sehr viele Daten benötigt, sehr fehleranfällig
-> Unter Grundverdacht: Vorwurf = schwere Straftat, massenweise Daten zur weiteren Überprüfung
- Altersverifikation: Perso != Privatheit, Bilderkennung kann falsch sein
- Diskussion: Diskussionsfragen
- Zusammenfassung \& Ausblick: EU-Ding verschoben, weitere Vorschläge zur Umsetzung
- Literatur
